\documentclass[UTF8,11pt,fontset=none]{ctexart}
\usepackage[a4paper,margin=23mm,headheight=15pt,footskip=11mm]{geometry}
\usepackage{fontspec}
\IfFontExistsTF{TeX Gyre Termes}{\setmainfont{TeX Gyre Termes}}{\setmainfont{lmroman10-regular.otf}[BoldFont=lmroman10-bold.otf,ItalicFont=lmroman10-italic.otf,BoldItalicFont=lmroman10-bolditalic.otf]}
\IfFontExistsTF{TeX Gyre Heros}{\setsansfont{TeX Gyre Heros}}{\setsansfont{lmsans10-regular.otf}[BoldFont=lmsans10-bold.otf,ItalicFont=lmsans10-oblique.otf]}
\IfFontExistsTF{DejaVu Sans Mono}{\setmonofont{DejaVu Sans Mono}[Scale=MatchLowercase]}{\setmonofont{lmmono10-regular.otf}}
\IfFontExistsTF{Noto Serif CJK SC}{\setCJKmainfont{Noto Serif CJK SC}[AutoFakeSlant=0.15]}{\setCJKmainfont{FandolSong-Regular.otf}[BoldFont=FandolSong-Bold.otf,ItalicFont=FandolKai-Regular.otf]}
\IfFontExistsTF{Noto Sans CJK SC}{\setCJKsansfont{Noto Sans CJK SC}}{\setCJKsansfont{FandolHei-Regular.otf}}
\IfFontExistsTF{Noto Sans Mono CJK SC}{\setCJKmonofont{Noto Sans Mono CJK SC}}{\setCJKmonofont{FandolFang-Regular.otf}}
\usepackage{amsmath,amssymb,amsthm,mathtools,mathrsfs,esint,bm}
\usepackage{booktabs,longtable,array,calc}
\usepackage{enumitem}
\setlist{nosep,leftmargin=1.7em,topsep=3pt}
\usepackage{fvextra}
\DefineVerbatimEnvironment{Highlighting}{Verbatim}{breaklines,breakanywhere,fontsize=\small}
\usepackage{xurl}
\usepackage[unicode,hidelinks,bookmarksnumbered]{hyperref}
\usepackage{bookmark}
\usepackage{fancyhdr}
\pagestyle{fancy}\fancyhf{}
\fancyhead[L]{\small 第一科学主线：近期研究全稿}
\fancyhead[R]{\small 2026-10-07}\fancyfoot[C]{\thepage}
\renewcommand{\headrulewidth}{0.3pt}
\providecommand{\tightlist}{\setlength{\itemsep}{1pt}\setlength{\parskip}{0pt}}
\providecommand{\passthrough}[1]{#1}
\newcommand{\Ric}{\operatorname{Ric}}
\newcommand{\Vol}{\operatorname{Vol}}
\newcommand{\Rm}{\operatorname{Rm}}
\newcommand{\supp}{\operatorname{supp}}
\newcommand{\Var}{\operatorname{Var}}
\newcommand{\tr}{\operatorname{tr}}
\newcommand{\avint}{\fint}
\newcommand{\Jcal}{\mathcal J}
\newcommand{\Ecal}{\mathcal E}
\newcommand{\dd}{\,\mathrm d}
\newcommand{\Rea}{\operatorname{Re}}
\setlength{\parindent}{2em}
\setlength{\parskip}{3pt plus 1pt}
\setlength{\emergencystretch}{3em}
\linespread{1.10}\allowdisplaybreaks[2]
\setlength{\abovedisplayskip}{8pt plus 2pt minus 2pt}
\setlength{\belowdisplayskip}{8pt plus 2pt minus 2pt}
\renewcommand{\arraystretch}{1.12}
\setcounter{tocdepth}{1}\setcounter{secnumdepth}{2}
\pdfstringdefDisableCommands{\def\fint{int}\def\mathcal#1{#1}\def\mathfrak#1{#1}\def\mathscr#1{#1}\def\operatorname#1{#1}\def\varepsilon{epsilon}\def\chi{chi}\def\beta{beta}\def\rho{rho}\def\tau{tau}\def\Phi{Phi}\def\sigma{sigma}\def\kappa{kappa}\def\nabla{D}\def\mathbb#1{#1}\def\({}\def\){}\def\textsuperscript#1{#1}\def\textsubscript#1{#1}}
\title{双取向实际类、统一间隙与 Kodaira 斜率}
\author{第一科学主线·近期研究分卷}
\date{2026年10月7日}
\begin{document}
\maketitle
\noindent\textbf{原稿编号：}R10。

本章以有限密度端点在固定值 \(B=2\) 的陈述为前置，联合同一度量上两个真实辛规范类的平方与交叉积，推导曲面丛的统一 Einstein 间隙。使用两个取向的规范理论本身已有经典先例；这里整理的是定量统一宽度的推导。\cite{R10,FD,Kotschick,LeBrunKodaira}

\section{条件性主结论}\label{ch10:ux6761ux4ef6ux6027ux4e3bux7ed3ux8bba}

采用下文 FD2 前置，对所有闭定向曲面丛 \[
\Sigma_f\longrightarrow M\longrightarrow\Sigma_b,\qquad f,b\ge2,
\] 存在与两个亏格、单值、总空间和度量无关的 \(\varepsilon_{\rm SB}>0\)，使 \[
\mathcal J(g)+\left(1-\frac{3|\tau(M)|}{\chi(M)}\right)_+
\ge\varepsilon_{\rm SB}.\tag{10.1}\label{eq:r10:1}
\] Einstein 度量因此满足 \[
3|\tau(M)|\le(1-\varepsilon_{\rm SB})\chi(M).\tag{10.2}\label{eq:r10:2}
\] 经典双取向辛构造和严格不等式是输入；本章新增推导是与整个范畴统一的正宽度。\cite{Kotschick,LeBrunKodaira}

\section{唯一非经典解析前置 FD2}\label{ch10:ux552fux4e00ux975eux7ecfux5178ux89e3ux6790ux524dux7f6e-fd2}

使用项目 \path{O10_All_Finite_Density.md} 第1节及第4节所陈述的定理，在单个固定密度 B=2 上：

\begin{quote}
\textbf{FD2.} 在原实际共形单极子供给下，若同一实际类 c 满足 \[
p_c=(c^+)^2\ge A_c>0,\quad c^2\ge T_+=2\chi+3\tau,
\quad\chi/A_c\le2,
\] 则存在 \(c_{\mathrm{FD}}(2)>0\)，对每一个允许度量及全部 \(\beta\ge2\) 有 \[
\fint|\operatorname{Ric}+g|^2+\frac{D_\beta(g,c)}V
\ge c_{\mathrm{FD}}(2)\frac{(\beta-2)^2}{\beta^2},
\] 其中 \[
D_\beta(g,c)=(1+\beta^{-1})\int s^2
+6(\beta+1)\int|W^+|^2-72\pi^2p_c.
\] 常数不依赖几何紧性、直径、非塌缩、Sobolev 常数或具体流形。
\end{quote}

本文仅在反取向 \(\bar M\) 的 \(s=-4\) 代表上使用 FD2。固定

\[
c=\min\{1,c_{\mathrm{FD}}(2)\},\qquad
\boxed{\varepsilon_{\mathrm{SB}}=(c/200)^2.}
\tag{10.3}\label{eq:r10:3}
\]

取常数200以简化估计；不求其最佳值。

经典前置为：Thurston 曲面丛辛形式构造；Taubes 的辛典范类 SW 非消失以及 \(b^+>1\) 时实际单极子供给；SW 标量预算；Yamabe 归一化；Chern--Weil；Hodge 极化。项目 N1 继承稿已记录所需的共形比较，下面亦给出计算。

\section{曲面丛提供的两个真实类}\label{ch10:ux66f2ux9762ux4e1bux63d0ux4f9bux7684ux4e24ux4e2aux771fux5b9eux7c7b}

写 \(e=\chi(M)=4(f-1)(b-1)>0\)。必要时反转基底取向，从而反转总空间取向，使 \(\tau\ge0\)。\(e,\mathcal J,|\tau|\) 不受这个选择影响。

取实定向平面丛分裂

\[
TM\simeq T_{\mathrm{vert}}M\oplus\pi^*T\Sigma_b.
\]

令 \(x=e(\pi^*T\Sigma_b)\)、\(y=e(T_{\mathrm{vert}}M)\)。在所选总空间取向上，

\[
x^2=0,\qquad x\cdot y=e,\qquad y^2=3\tau.
\tag{10.4}\label{eq:r10:4}
\]

第一式因 x 来自二维基底；第二式为 Euler 类乘积；第三式来自 \(p_1(TM)=x^2+y^2\) 与签名公式。

Thurston 构造分别对基底面积项取正、负符号，得到两个总空间取向上的辛形式。其相容几乎复结构分别与水平/垂直块的取向选择同伦。因此可取两个实际典范或反典范类

\[
a=x+y,\qquad b=x-y.
\]

b 与反取向结构的实际典范类只可能相差整体符号，该符号不影响下文平方和绝对配对。于是

\[
\boxed{
 a^2=2e+3\tau,\quad b^2[M]=-2e+3\tau,
 \quad a\cdot b=-3\tau.
}
\tag{10.5}\label{eq:r10:5}
\]

b 是 \(\bar M\) 上的真实辛典范类，不是从一个形式交叉向量猜出的单极子类。

由 \(\pi_1(M)\twoheadrightarrow\pi_1(\Sigma_b)\)，有 \(b_1(M)\ge2b\ge4\)。当前取向满足

\[
b_2^+(M)=\frac{e-2+2b_1+\tau}{2}>1.
\]

因此 a 具有全度量非消失供给。其正平方 \(a^2=2e+3\tau>0\) 排除了非负标量曲率的 Yamabe 代表：该代表上的未扰动 SW 方程与 Weitzenböck 恒等式会迫使旋量为零、\(a^+=0\)，与 \((a^+)^2\ge a^2>0\) 矛盾。

故 M 的每个共形类均为负 Yamabe 类，所有度量均有 R\textgreater0。本文不对 Ricci-flat 度量使用未定义的 \(4P/R\)。

反取向的 \(b_2^+\) 条件将在下文以正取向曲率预算推出，避免先假定一个待证的签名上界。

\section{转到同一个负 Yamabe 代表}\label{ch10:ux8f6cux5230ux540cux4e00ux4e2aux8d1f-yamabe-ux4ee3ux8868}

若原度量 \(\mathcal J(g)\ge1\)，\eqref{eq:r10:1}直接成立。故只考虑 \(\mathcal J(g)<1\)。

令 \(\hat g\in[g]\) 为 \(s_{\hat g}=-4\) 的代表。Weyl 能量的共形不变性与 Chern--Weil 给出

\[
P(g)-P(\hat g)=\Delta S/12,\quad
R(g)-R(\hat g)=\Delta S/3,
\quad\Delta S=\int_gs^2-\int_{\hat g}s^2\ge0.
\]

因而 \(\mathcal J(g)\) 是 \(\mathcal J(\hat g)\) 与1的加权平均。原 \(\mathcal J(g)<1\) 迫使

\[
j:=\mathcal J(\hat g)\le\mathcal J(g)<1.
\tag{10.6}\label{eq:r10:6}
\]

以下全部度量、Hodge 投影和实际解都使用同一个 \(\hat g\)。记

\[
t=3\tau/e\ge0,\quad u=V/(2\pi^2e),\quad
r=P/V=\frac{4j}{4-j},\quad
w=\fint|W^-|^2.
\tag{10.7}\label{eq:r10:7}
\]

由于 \(\hat g\) 为常标量代表，原归一化缺陷还满足 \(\mathcal E(\hat g)=j\)，但证明不需要从这个等式推出任何极小值取得性。

\section{先仅使用正取向实际类和两手性 Chern--Weil}\label{ch10:ux5148ux4ec5ux4f7fux7528ux6b63ux53d6ux5411ux5b9eux9645ux7c7bux548cux4e24ux624bux6027-chernweil}

设 \(U=(a^+)^2/e\)。实际 SW 标量预算

\[
16V\ge32\pi^2(a^+)^2
\]

及交叉平方给出

\[
2+t\le U\le u.
\tag{10.8}\label{eq:r10:8}
\]

反手性 Chern--Weil 公式给出精确关系

\[
\boxed{w=\frac{2-t}{u}-\frac13+\frac{j}{4-j}\ge0.}
\tag{10.9}\label{eq:r10:9}
\]

由于 j\textless1，\eqref{eq:r10:9}首先推出 t\textless2，并给出

\[
2+t\le u\le C(j)(2-t),\qquad
C(j)=\frac{3(4-j)}{4(1-j)}.
\tag{10.10}\label{eq:r10:10}
\]

于是

\[
(16-7j)t\le16+2j,
\quad
\boxed{t\le1+\frac{9j}{16-7j}\le1+j.}
\tag{10.11}\label{eq:r10:11}
\]

现在 t\textless2 推出

\[
b_2^-(M)=\frac{e-2+2b_1-\tau}{2}
>\frac e6+b_1-1>1.
\]

所以反取向的辛典范类 b 也具有全度量实际供给，且

\[
b^2[\bar M]=e(2-t)>0.
\]

由其实际标量预算，令 \(Z=\|b^-\|^2/e\)（负 Hodge 范数取正值），则

\[
2-t\le Z\le u.
\tag{10.12}\label{eq:r10:12}
\]

这里不存在从同一自旋丛中任意改取向的步骤：a 和 b 分别由两个真实辛结构提供。

\section{同一个度量上的交叉配对补出关键极化信息}\label{ch10:ux540cux4e00ux4e2aux5ea6ux91cfux4e0aux7684ux4ea4ux53c9ux914dux5bf9ux8865ux51faux5173ux952eux6781ux5316ux4fe1ux606f}

置

\[
d=j+(1-t)_+.
\]

如果 \(d\ge1/256\)，则因 \(\varepsilon_{\mathrm{SB}}\le1/40000\)，结论已成立。暂设

\[
0<d\le1/256.
\tag{10.13}\label{eq:r10:13}
\]

设

\[
X=U-(2+t)=\|a^-\|^2/e\ge0,
\qquad Y=Z-(2-t)=\|b^+\|^2/e\ge0.
\]

由\eqref{eq:r10:10}--\eqref{eq:r10:11}可得

\[
1-d\le t\le1+d,\qquad 3-d\le u\le3+5d<4,
\] \[
\boxed{X\le4d,\quad w\le d,\quad r\le4d/3.}
\tag{10.14}\label{eq:r10:14}
\]

给出常数核对：当 j\(\le\)1/4 时，\(C(j)\le3+3j\)。记 \(d_-=(1-t)_+\)，则

\[
X\le4d_-+3j(1+d_-)\le4(j+d_-)=4d.
\]

\eqref{eq:r10:9}和 \(u\ge2+t\) 给

\[
w\le\frac{4(1-t)}{3(2+t)}+\frac{j}{4-j}
\le\frac12d_-+\frac13j\le d.
\]

现在使用独立于两项平方预算的真实拓扑配对。Hodge 正交分解和 Cauchy--Schwarz 使\eqref{eq:r10:5}变成

\[
\boxed{t\le\sqrt{UY}+\sqrt{XZ}.}
\tag{10.15}\label{eq:r10:15}
\]

在 \(X\le4d,\ Z\le u<4\) 下，有 \(\sqrt{XZ}\le4\sqrt d\)。由于 \(t-4\sqrt d>0\)，

\[
\frac Zu\ge\frac{2-t}{u}
+\frac{(t-4\sqrt d)^2}{u^2}.
\tag{10.16}\label{eq:r10:16}
\]

第一项由 \(w\ge0\) 控制：

\[
\frac{2-t}{u}\ge\frac13-\frac j{4-j}\ge\frac13-\frac d3.
\]

第二项满足

\[
\frac{(t-4\sqrt d)^2}{u^2}
\ge\frac{(1-d-4\sqrt d)^2}{(3+5d)^2}
\ge\frac19-\frac89\sqrt d-\frac{16}{27}d.
\]

最后一个不等式由

\[
(1-d-4\sqrt d)^2\ge1-2d-8\sqrt d>0,
\quad
(3+5d)^{-2}\ge\frac19-\frac{10}{27}d>0
\]

逐项相乘得到。因此

\[
\frac Zu\ge\frac49-\frac89\sqrt d-\frac{25}{27}d
\ge\boxed{\frac49-2\sqrt d.}
\tag{10.17}\label{eq:r10:17}
\]

这是新接通的关键：若正取向接近复双曲拓扑边界且无迹 Ricci 很小，正取向实际类几乎完全自对偶；它与反取向实际类的固定交叉积，迫使后者的正投影（在反取向下）进入原禁止端点的数值范围。单独知道两个类的平方无法得到\eqref{eq:r10:15}。

\section{反取向端点矛盾；常数对全部曲面丛统一}\label{ch10:ux53cdux53d6ux5411ux7aefux70b9ux77dbux76feux5e38ux6570ux5bf9ux5168ux90e8ux66f2ux9762ux4e1bux7edfux4e00}

在 \(\bar M\) 上使用 b，选取预算

\[
A_-=e(2-t)=b^2[\bar M]=T_+(\bar M)>0.
\]

投影 \(p_b=Ze\ge A_-\)，同一类平方条件恰取等号。而且

\[
\frac{\chi(\bar M)}{A_-}=\frac1{2-t}
\le\frac1{1-d}<2.
\tag{10.18}\label{eq:r10:18}
\]

所以无论两个亏格和单值怎样变化，FD2都在\textbf{同一个固定密度2}上使用。

反取向参数缺陷为

\[
\bar\kappa_\beta
=r+16(1+\beta^{-1})+6(\beta+1)w-36\frac Zu.
\tag{10.19}\label{eq:r10:19}
\]

取有限参数 \(\beta=d^{-1/2}\ge16\)。由\eqref{eq:r10:14}、\eqref{eq:r10:17}，

\[
\bar\kappa_\beta
\le\frac43d+16\sqrt d+6(\sqrt d+d)+72\sqrt d
=94\sqrt d+\frac{22}{3}d
<100\sqrt d.
\tag{10.20}\label{eq:r10:20}
\]

FD2同时给出

\[
\bar\kappa_\beta\ge c(1-2\sqrt d)^2
\ge\frac{49}{64}c\ge\frac12c.
\tag{10.21}\label{eq:r10:21}
\]

于是 \(d\ge(c/200)^2\)。

还须处理 d=0，而不能除以零：此时 j=0,t=1,u=3,X=0,w=0；\eqref{eq:r10:15}直接给 \(Z\ge4/3\)。因此对每个有限 \(\beta\ge2\)，

\[
\bar\kappa_\beta\le16/\beta.
\]

这与FD2在 \(\beta\to\infty\) 时的正下界矛盾。没有把 \(\beta=\infty\) 当成一个原定理参数。

所有分支合并得到

\[
j+(1-t)_+\ge(c/200)^2.
\]

再用 \(j\le\mathcal J(g)\) 和取向可反转，证明\eqref{eq:r10:1}。证毕。

\section{Einstein、Kodaira与全度量含义}\label{ch10:einsteinkodairaux4e0eux5168ux5ea6ux91cfux542bux4e49}

\subsection{1 Einstein与非Einstein度量}\label{ch10:einsteinux4e0eux975eeinsteinux5ea6ux91cf}

实际正预算排除了非负标量曲率，因此这里的Einstein度量必具有负Einstein常数，\(R>0\)、\(\mathcal J=0\)。将其代入\eqref{eq:r10:1}，得到\eqref{eq:r10:2}。

对固定流形，若 \(d_{\mathrm{top}}=(1-3|\tau|/e)_+<\varepsilon_{\mathrm{SB}}\)，则所有度量都有

\[
\mathcal J(g)\ge\varepsilon_{\mathrm{SB}}-d_{\mathrm{top}}>0.
\]

利用正取向实际标量预算还给出

\[
P(g)\ge2\pi^2 e(2+t)
\bigl(\varepsilon_{\mathrm{SB}}-d_{\mathrm{top}}\bigr).
\]

这不是仅判断一个局部修复是否成功，而是约束所有光滑竞争度量。

\subsection{2 Kodaira纤维化}\label{ch10:kodairaux7ea4ux7ef4ux5316}

对紧复曲面的光滑Kodaira纤维化，基底和纤维亏格均至少2，总空间无有理曲线。它是极小一般型曲面；其规范丛为ample，因而具有负Kähler--Einstein度量。这里使用经典曲面论与Aubin--Yau存在定理，不将其作为新证明。

在复取向，\(c_1^2=2e+3\tau\)、\(c_2=e\)。由\eqref{eq:r10:2}得到

\[
\boxed{c_1^2(M)\le(3-\varepsilon_{\mathrm{SB}})c_2(M).}
\tag{10.22}\label{eq:r10:22}
\]

也可仅把\eqref{eq:r10:22}陈述为对已知承载负KE度量的Kodaira纤维化的推论，从而不需在本稿中展开规范丛amplicity的经典证明。

其意义是把经典严格 \(c_1^2<3c_2\) 改为独立于亏格和单值的正宽度。本文没有证明最佳宽度，更没有证明 \(8/3\) 是Kodaira Chern斜率的最优上界。

\subsection{3 一个精确校准}\label{ch10:ux4e00ux4e2aux7cbeux786eux6821ux51c6}

在负KE的Kodaira情形，记 \(r_K=c_1^2/e\)。其反取向典范类是原复结构的(1,1)类，故

\[
p_b=\frac{4e^2}{c_1^2},\quad
V=2\pi^2 c_1^2,\quad
\fint|W^-|^2=\frac{4(3-r_K)}{3r_K}.
\]

由此\eqref{eq:r10:19}化为

\[
\bar\kappa_\beta
=16(1+\beta^{-1})
+8(\beta+1)\frac{3-r_K}{r_K}
-\frac{144}{r_K^2}.
\tag{10.23}\label{eq:r10:23}
\]

当 \(r_K=2,\beta=2\) 时，右侧为0，符合允许的乘积KE端；当形式上 \(r_K=3\) 时，它等于 \(16/\beta\)，恰好进入FD2排除的极限。这只是符号与常数校准，不假定存在斜率3的光滑Kodaira纤维化。

\section{与公开斜率问题的关系}\label{ch10:ux4e0eux516cux5f00ux659cux7387ux95eeux9898ux7684ux5173ux7cfb}

此前整理的 Catanese 综述第1.3节记录 LeBrun 的问题：是否存在对全部 Kodaira 纤维化统一的 \(\epsilon>0\)，使 \(c_1^2/c_2\le3-\epsilon\)；Polizzi 的 Remark 4.10 以斜率能否任意逼近 \(3\) 表述同一边界问题。\cite{Catanese2016,Polizzi2022}

本章在 FD2 前置下导出的 \(c_1^2/c_2\le3-\varepsilon_{\rm SB}\) 排除这种逼近。该陈述不确定最佳斜率、不证明 \(8/3\) 为上确界；其依赖是 FD2、两个真实辛规范类和同度量交叉配对。

\begin{thebibliography}{P99}
\small
\bibitem[R10]{R10} 本会话研究原稿，\emph{双取向实际类、统一间隙与 Kodaira 斜率}，2026-10-07。\path{Line1_Two_Orientation_Global_Gap.md}。

\bibitem[P05]{FD} 所提供网页续研，\emph{O10 closes at every fixed finite Euler density}，2026-09-29，\path{O10_All_Finite_Density.md}。

\bibitem[E01]{Kotschick} D. Kotschick, \emph{On regularly fibered complex surfaces}, Geometry \& Topology Monographs \textbf{2} (1999), 291--298, Theorem 3. \url{https://arxiv.org/abs/math/9911255}.

\bibitem[E02]{LeBrunKodaira} C. LeBrun, \emph{Diffeomorphisms, symplectic forms and Kodaira fibrations}, Geometry \& Topology \textbf{4} (2000), 451--456. \url{https://arxiv.org/abs/math/0005195}.

\bibitem[E05]{Catanese2016} F. Catanese, \emph{Kodaira fibrations and beyond: methods for moduli theory}, arXiv:1611.06617, §1.3. \url{https://arxiv.org/abs/1611.06617}.

\bibitem[E06]{Polizzi2022} F. Polizzi, \emph{Diagonal double Kodaira structures on finite groups}, arXiv:2002.01363v2, Remark 4.10. \url{https://arxiv.org/abs/2002.01363}.

\end{thebibliography}
\end{document}
